\section{Two-sided ideal receipts}
\label{sec:ideal}
\status{Producer, independent Python replay, exact oracle comparisons, and
adversarial tests are implemented. Lean source reification is not implemented.}

\subsection{The receipt and its soundness}

For supplied relations $r_1,\ldots,r_s$ and target $p$, a receipt consists of a
finite list of rational coefficients, relation indices, and two words:
\begin{equation}\label{eq:ideal}
 p=\sum_{\ell=1}^{t} c_\ell u_\ell r_{j_\ell}v_\ell.
\end{equation}
The left and right words are separately represented fields. The original
relations and target are arguments to replay, not assertions made by the receipt.

\begin{theorem}[Two-sided receipt soundness]\label{thm:ideal-sound}
If identity~\eqref{eq:ideal} holds in $\F$, then under every evaluation into an
associative unital $\Q$-algebra satisfying $r_j=0$, the evaluated target is zero.
\end{theorem}
\begin{proof}
Evaluation is an algebra homomorphism, so it preserves the finite sum and each
product in~\eqref{eq:ideal}. Every summand has the factor $r_{j_\ell}=0$.
No commutation, division by an operator, or cancellation of a nonzero element is
used. Rational coefficients act through the central scalar map.
\end{proof}

A receipt need not mention a Gr\"obner basis, a monomial ordering, or a minimal
derivation. Those are possible discovery mechanisms, not facts needed for this
conclusion. Compared with the existing commutative multiplier format, the
essential change is not replacing an implementation name: it is preserving both
sides of each relation occurrence.

\subsection{A concrete bounded producer}

Choose a degree bound $D\ge\deg p$. For each relation of degree $d_j$, enumerate
all pairs of words $(u,v)$ with $|u|+|v|\le D-d_j$ and expand the column
$u r_jv$. A sparse exact elimination basis tracks each admitted row as a rational
combination of these original columns. Reducing $p$ against that basis either
returns the coefficients in~\eqref{eq:ideal} or a nonzero remainder.

\begin{lstlisting}
columns := [(j,u,v, expand(u * relation[j] * v))
            for j,u,v with len(u)+degree(relation[j])+len(v) <= D]
selected := incidenceComponent(target.support, columns)
basis := empty sparse exact echelon basis with provenance
for column in selected:
    insert(column, original_column_index)
(coefficients, remainder) := reduce(target, basis)
if remainder != 0:
    return UNKNOWN_OUTSIDE_BOUNDED_CONTEXT_SPAN
receipt := original-context terms selected by coefficients
return receipt only after independent expansion succeeds
\end{lstlisting}

The prototype performs Gaussian elimination over \code{fractions.Fraction}.
A production backend can replace it by sparse fraction-free or modular exact
linear algebra without changing the receipt format. The relevant finite
feasibility problem is linear; there are no unknown generators being multiplied
by unknown multipliers.

\paragraph{Complexity.}
There are at most
\begin{equation}\label{eq:counts}
 R_D=\sum_{k=0}^{D}m^k\quad\hbox{word rows},\qquad
 C_D=\sum_{j=1}^{s}\sum_{\ell=0}^{D-d_j}(\ell+1)m^\ell
 \quad\hbox{context columns}.
\end{equation}
An upper sum with $D-d_j<0$ is empty. The factor $\ell+1$ counts the possible
split between the two contexts of a length-$\ell$ word. Dense elimination costs
$O(R_D C_D\min(R_D,C_D))$ field operations; sparse behavior and rational bit
growth determine the practical cost. These counts are exponential in $D$, not
a claim of a generally efficient noncommutative theorem prover. The implemented
column cap is 20,000; exceeding it raises an explicit budget exception that a
controller must map to \unknown{}.

\subsection{Exact incidence slicing before elimination}

Construct a bipartite graph with word rows on one side and context columns on
the other; an edge means the coefficient of that word in that column is nonzero.
Keep every connected component that meets the target's support. This is more
precise than merely keeping all columns mentioning one of the target's variables.

\begin{proposition}[Component slicing is exact]\label{prop:slicing}
The target lies in the span of all columns if and only if it lies in the span of
the retained columns.
\end{proposition}
\begin{proof}
One implication is immediate. For the other, partition a hypothetical solution
into graph components. Distinct components have disjoint row supports. Every
component disjoint from the target must sum to zero on its own rows, independently
of the components meeting the target. Removing all such contributions leaves a
solution for the target using only retained columns.
\end{proof}

No heuristic premise deletion is hidden in this argument. Conversely, slicing
only one hop from the target would be wrong: a column may introduce another word
whose cancellation requires a further column. The prototype uses a worklist
until the whole incidence component is exhausted.

The executed implementation assembles \emph{all} columns before constructing
the graph. Its measured benefit is reduced elimination, not avoidance of the
exponential column-assembly cost. A future lazy context index may improve that
cost, but it needs a completeness argument for enumerating every incident
column. The current result must not be described as such a lazy algorithm.

\subsection{Why an unsuccessful bounded search is not a refutation}

Set
\[
 r_U=U(1-AB)-1,\qquad r_V=(1-BA)V-1,\qquad p=UA-AV.
\]
The following exact identity holds in the free algebra:
\begin{equation}\label{eq:push}
 p=r_UAV-UA r_V.
\end{equation}
The target has degree two, but the displayed context products have degree five.
In the delivered experiment, context bounds two, three, and four all return
\unknown{}, whereas bound five returns a receipt. This is a direct operational
control against equating ``target degree'' with a sufficient proof degree for
inhomogeneous presentations.

There is a sharper semantic warning. The relation $xy=1$ does not imply $yx=1$
in every algebra: left and right shifts on the countably generated vector space
provide a counterexample. If $S e_i=e_{i+1}$ and $L e_0=0$, $L e_{i+1}=e_i$,
then $LS=1$ but $SL\ne1$. Yet square matrices over a field with $XY=I$ also
satisfy $YX=I$, since a right inverse in finite dimension is an inverse. Thus a
finite-matrix search is not complete for arbitrary inhomogeneous algebra
counterexamples. The homogeneous restriction in the next section is doing
substantial work.

\section{A complete homogeneous fragment with matrix witnesses}
\label{sec:homogeneous}

Let $\F_k$ be the rational span of words of length $k$. A relation is homogeneous
when all its nonzero monomials have the same degree. Different relations may
have different degrees; zero relations are harmless. Let $I$ be their two-sided
ideal and set
\[
 \F_{\le D}=\bigoplus_{k=0}^{D}\F_k,\qquad
 \F_{>D}=\bigoplus_{k>D}\F_k.
\]
The second space is a two-sided ideal: multiplying a word of length greater
than $D$ by any word cannot shorten it.

\subsection{Why degree \texorpdfstring{$D$}{D} suffices}

\begin{lemma}[Degree-bounded membership]\label{lem:graded-membership}
For homogeneous relations and $p\in\F_{\le D}$, membership $p\in I$ is equivalent
to membership in the finite span
\[
 \Span_\Q\{u r_jv:\ |u|+d_j+|v|\le D\}.
\]
\end{lemma}
\begin{proof}
The reverse implication follows from the definition of an ideal. For the forward
implication, write a finite ideal representation of $p$ and expand both
polynomial multipliers into words. Each resulting context product is
homogeneous. Group these products by total degree. Every group of degree greater
than $D$ sums to zero, since $p$ has no component in that degree. Delete those
groups. The remaining products satisfy the asserted degree bound.
\end{proof}

This proof does not appeal to termination of noncommutative completion or to a
Noetherian property. It isolates a finite-dimensional vector-space problem at
the exact required degree. It needs no global completion result for the presented algebra.

\begin{theorem}[Exact graded alternative]\label{thm:graded-alternative}
Given finitely many homogeneous rational relations and a target of degree at
most $D$, exact elimination with all contexts of total degree at most $D$
returns either a receipt~\eqref{eq:ideal}, or rational square matrices
$X_0,\ldots,X_{m-1}$ and a rational vector $v$ such that
\[
 r_j(X_0,\ldots,X_{m-1})=0\quad\hbox{for every }j,
 \qquad p(X_0,\ldots,X_{m-1})v\ne0.
\]
A witness dimension no larger than $R_D$ in~\eqref{eq:counts} suffices.
\end{theorem}
\begin{proof}
If elimination succeeds, Theorem~\ref{thm:ideal-sound} applies. Otherwise
Lemma~\ref{lem:graded-membership} gives $p\notin I$. Because $I$ is graded,
\[
 p\notin I+\F_{>D}.
\]
Indeed, the degree-at-most-$D$ part of any element of $I$ is again in $I$.
Consider the finite-dimensional quotient algebra
\[
 A_D=\F/(I+\F_{>D}).
\]
Its dimension is at most $R_D$. The image of $p$ is nonzero, so this quotient is
not the zero algebra and its identity is nonzero. Let $X_i$ be the matrix of
left multiplication by the image of $x_i$ in any rational basis of $A_D$, and
let $v$ represent its identity. Left multiplication is a unital algebra
homomorphism. Consequently all input relations evaluate to the zero matrix,
whereas $p(X)v$ is exactly the nonzero coordinate vector of the image of $p$.
\end{proof}

\subsection{An executable construction, not an existence-only argument}

The producer forms the full degree-bounded echelon basis. Nonpivot words form a
basis of $A_D$. To construct column $j$ of $X_i$, prepend $x_i$ to the $j$th
basis word and reduce it; a word longer than $D$ reduces to zero by truncation.
Reduce the empty word to construct $v$.

\begin{lstlisting}
E := exact basis of all relation contexts through degree D
B := all nonpivot words of length at most D
for atom i:
    for basis word B[j]:
        X[i][:,j] := coordinates(normal_form(i :: B[j]))
v := coordinates(normal_form(empty_word))
return X,v only after matrix replay on the ORIGINAL relations
\end{lstlisting}

Unlike the equality receipt, this producer must use the \emph{full} relation
space, not just the target's incidence component. A sliced equality proof is
sound; a quotient that omits unrelated relations need not model the whole
presentation. The implementation uses full elimination here deliberately.

The checker does not reconstruct a quotient, inspect pivots, verify a basis, or
trust the producer's homogeneity assertion. It multiplies the supplied rational
matrices and verifies every complete relation matrix, then the separating target
vector. An accidentally inapplicable construction cannot pass merely because it
looks like a quotient. The producer rejects inhomogeneous input before attempting
this construction; the matrix checker itself can accept a directly valid
countermodel for any equality-schema relation list.

\paragraph{A crucial independent test.}
Take the purported relation $X=0$, the matrix $X=\operatorname{diag}(0,1)$,
target $1$, and vector $e_1$. Then $Xe_1=0$ and $1e_1\ne0$, but $X\ne0$.
Checking the relation only on the witness vector would accept a false
countermodel. The supplied unit suite rejects exactly this object.

\subsection{What the completeness statement does and does not imply}

For homogeneous presentations, the theorem decides the universal algebra
consequence and, equivalently, the consequence over rational matrices of all
finite sizes. It does not decide the corresponding consequence in one selected
algebra, nor with omitted positivity, involution, or invertibility assumptions.
The quotient matrices need not be symmetric, and their dimension is whatever
the quotient needs. They are not automatic counterexamples to a goal about
$2\times2$ Hermitian matrices.

Operational caps are a separate boundary. The implemented producer refuses an
ambient word count above 4,000 or a quotient dimension above 100. Such a refusal
is \unknown{} under a resource limit, even though the uncapped mathematical
procedure terminates. No matrix witness is promised below an arbitrarily chosen
user cap.

\subsection{The algorithmic consequence for Forge}

This worker has an unusually clean escalation policy. Homogeneous equalities
need only one degree level, followed by either identity replay or matrix replay.
Inhomogeneous equalities use increasing context bounds and never turn a
nonzero bounded remainder into a source-level refutation. This distinction is
available syntactically before the expensive search begins. It is more useful
than a generic instruction to ``try a bigger algebra solver.''
